% !TEX root = ./main.tex
%%%%%%%%%%%%%%%%%%%%%%%%%%%%%%%%%%%%%%%%%%%%%%%%%%%%%%%%%%%%%%%%%%%%%%%%%%%%%%%%%%%%%%%%%%
% Autor: Damir Agačević, 2025.
% damir@agacevic.com
% https://github.com/theDDA/
% GNU GPLv3
%%%%%%%%%%%%%%%%%%%%%%%%%%%%%%%%%%%%%%%%%%%%%%%%%%%%%%%%%%%%%%%%%%%%%%%%%%%%%%%%%%%%%%%%%%

% A4, 12pt font
\documentclass[12pt,a4paper,final,oneside]{memoir}

% Mogu se direktno unositi čćžđš.
\usepackage[utf8]{inputenc}
% Pdf daje npr. š umjesto ˇs pri copy/paste.
\usepackage{cmap}
\usepackage[T1]{fontenc}

% Lokalizacija
% TeX Live nema bosanska pravila za prelamanje riječi, pa se koriste srpska.
\usepackage{babel}
\babelprovide[import,main,hyphenrules=serbian]{bosnian}
\renewcommand\bosnianhyphenmins{22}
% Kao i za serbian: bez dodatnog razmaka nakon tačke i uvučen prvi pasus
% nakon naslova.
\addto\extrasbosnian{\frenchspacing}
\usepackage{indentfirst}

% Navodnici usklađeni s jezikom, traži ga biblatex uz babel.
\usepackage{csquotes}
\DeclareQuoteAlias{serbian}{bosnian}

% Proširenje matematike
\usepackage{amsmath, amssymb}
\usepackage{mathtools}

% Komanda \dd daje znak diferencijala koji je uspravan i malo odvojen
% od podintegralne funkcije, npr. $\int x \dd{x}$. Parcijalni izvod: \pdv{f}{x}.
\usepackage{physics}

% Feynmanova slash notacija, npr. $\slashed{p}$.
\usepackage{slashed}

% Font teksta i matematike prema izboru \stil u main.tex.
% Mora se učitati nakon amsmath.
% Svaki stil definiše i \fontnaslova, font naslova poglavlja i sekcija.
\ifcase\stil\relax
	\PackageError{preamble}{Nepoznat stil \stil}{U main.tex izabrati \string\stil{1}, {2} ili {3}.}
\or% !TEX root = ../main.tex
% Stil 1: Computer Modern.
\usepackage{lmodern}

\newcommand\fontnaslova{\sffamily}

\or% !TEX root = ../main.tex
% Stil 2: Helvetica (TeX Gyre Heros).
\usepackage{lmodern}
\usepackage[scale=0.95]{tgheros}
\renewcommand{\familydefault}{\sfdefault}

% Sans-serif matematika uz Heros. Mijenja \sfdefault, pa se Heros vraća odmah iza.
\usepackage{sansmathfonts}
\renewcommand{\sfdefault}{qhv}

\newcommand\fontnaslova{\sffamily}

\or% !TEX root = ../main.tex
% Stil 3: STIX.
% STIX nema sans-serif varijantu, pa su i naslovi serifni.
\usepackage{stix}
\usepackage[scaled=.85]{inconsolata}

\newcommand\fontnaslova{\rmfamily}

\else
	\PackageError{preamble}{Nepoznat stil \stil}{U main.tex izabrati \string\stil{1}, {2} ili {3}.}
\fi

% Uvlačenje paragrafa. Odkomentarisati da se prvi red paragrafa ne uvlači.
% \setlength{\parindent}{0pt}

% Slike
\usepackage{graphicx}
\graphicspath{{slike/}}

% Umetanje vanjskih PDF stranica, npr. nacrta u dodatku.
\usepackage{pdfpages}

% Crtanje dijagrama u tikz
\usepackage{tikz}
\usetikzlibrary{calc}

% Feynmanovi dijagrami
\usepackage[compat=1.1.0]{tikz-feynman}

% Grafici. Stil "grafik" daje jednak izgled svim graficima u radu:
% \begin{axis}[grafik, xlabel=..., ...]
\usepackage{pgfplots}
\usepgfplotslibrary{groupplots}
\pgfplotsset{
	compat=1.18,
	grafik/.style={
			scale only axis,
			width=0.72\linewidth,
			height=0.45\linewidth,
			tick align=outside,
			tick pos=left,
			every axis plot/.append style={line width=0.9pt},
			label style={font=\small},
			tick label style={font=\small},
			title style={font=\small},
			legend style={font=\small, draw=black, fill=white, cells={anchor=west}},
		},
}

% Proširenje za tabele
\usepackage{booktabs}
\usepackage{multirow}
\usepackage{makecell}
% Malo zraka između redova -- razlomci i indeksi inače dodiruju linije.
\renewcommand{\arraystretch}{1.15}
% Tekstualne kolone fiksne širine, lijevo poravnate (p kolone inače razvlače razmake).
\newcolumntype{L}[1]{>{\raggedright\arraybackslash}p{#1}}

% Proširenje za floats, potrebno za komandu [H] kojom se forsira tačno mjesto slike/tabele.
\usepackage{float}

% Rotirani floatovi, npr. \begin{sidewaystable}
\usepackage{rotating}

% Manji font opisa slika, opisi tabela iznad tabele, subcaption za opise u grupama slika.
\usepackage{caption}
\usepackage{subcaption}
\captionsetup{font=small,justification=justified,format=hang}
% Opisi tabela idu iznad tabele.
\captionsetup[table]{position=above}

% Lakše unošenje jedinica, npr. $g = \SI{9.81}{\meter\per\second\squared}$.
\usepackage{siunitx}
% physics i siunitx oba definišu \qty; ovdje \qty znači isto što i \SI.
\AtBeginDocument{\RenewCommandCopy\qty\SI}
% Ne razdvajati decimale u grupe po tri (0.17067, ne 0.170 67).
\sisetup{
	group-digits=none,
	input-decimal-markers={.,},
	range-phrase={ do },
	range-units=single,
}
% Odkomentarisati samo ako baš treba decimalni zarez umjesto tačke.
% \sisetup{output-decimal-marker={,}}
\DeclareSIUnit{\barn}{b}

% Liste: description labele u istoj familiji kao naslovi.
\usepackage{enumitem}
\setlist[description]{font=\fontnaslova\bfseries, style=nextline,
	leftmargin=1.5em, itemsep=0.4em}
% Liste listi se ne uvlače dodatno.
\setlist[itemize]{leftmargin=*}

% Biblatex za citate, numerisani redom citiranja.
\usepackage[sorting=none]{biblatex}
\addbibresource{literatura.bib}
% Biblatex nema bosanski, uzima srpski koji je ekavski.
\DeclareLanguageMapping{bosnian}{serbian}
\DefineBibliographyStrings{bosnian}{
	urlseen         = {posjećeno},
}
% Jedan izvor se ne prelama preko dvije stranice.
\appto\bibsetup{\interlinepenalty=10000\relax}

% Bolja tipografija i kerning
\usepackage[activate={true,nocompatibility},final,tracking=true,kerning=true,spacing=true,factor=1100,stretch=10,shrink=10]{microtype}
\SetProtrusion{encoding={*},family={*},series={*},size={6,7}}
{1={ ,750},2={ ,500},3={ ,500},4={ ,500},5={ ,500},
	6={ ,500},7={ ,600},8={ ,500},9={ ,500},0={ ,500}}

% Linkovi
\usepackage{hyperref}
\usepackage{memhfixc}

% Dozvoljava prelamanje align okruženja na novu stranicu
\allowdisplaybreaks

% Funkcije kako se pišu na našim prostorima.
% Češće koristimo \varphi
\newcommand\vp{\varphi}
% Hiperbolne trigonometrijske:
\DeclareMathOperator{\tgh}{th}
% Vektorske operacije:
\DeclareMathOperator{\rrot}{rot}
\DeclareMathOperator{\ddiv}{div}
\DeclareMathOperator{\ggrad}{grad}

% Komande za unos podataka u naslovnu stranicu
\makeatletter
\newcommand\naslov[1]{\gdef\@naslov{#1}}
\newcommand\nasloveng[1]{\gdef\@nasloveng{#1}}
\newcommand\podnaslov[1]{\gdef\@podnaslov{#1}}
\newcommand\podnasloveng[1]{\gdef\@podnasloveng{#1}}
\newcommand\smjer[1]{\gdef\@smjer{#1}}
\newcommand\smjereng[1]{\gdef\@smjereng{#1}}
\newcommand\student[1]{\gdef\@student{#1}}
\newcommand\profesor[1]{\gdef\@profesor{#1}}
\newcommand\datum[1]{\gdef\@datum{#1}}
\newcommand\datumeng[1]{\gdef\@datumeng{#1}}
\makeatother

%%%%%%%%%%%%%%%%%%%%%%%%%%%%%%%%%%%%%%%%%%%%%%%%%%%%%%%%%%%%%%%%%%%%%%%%%%%%%%%%%%%%%%%%%%
% Izgled stranice
%%%%%%%%%%%%%%%%%%%%%%%%%%%%%%%%%%%%%%%%%%%%%%%%%%%%%%%%%%%%%%%%%%%%%%%%%%%%%%%%%%%%%%%%%%

% Margine i blok teksta
\settypeblocksize{*}{32pc}{1.618}
\setbinding{0.5cm}
\setlrmargins{1.4in}{*}{*}
\setulmargins{*}{*}{1.3}

% Slova s kvačicom (npr. Č) su viša od običnog velikog slova, pa živa glava
% preraste \onelineskip. Malo više prostora za glavu.
\setheadfoot{15pt}{2\onelineskip}
\setheaderspaces{*}{2\onelineskip}{*}

\checkandfixthelayout

% Naslovi poglavlja i sekcija
\makechapterstyle{mychapterstyle}{%
	\renewcommand{\chapnamefont}{\LARGE\fontnaslova\bfseries}%
	\renewcommand{\chapnumfont}{\LARGE\fontnaslova\bfseries}%
	\renewcommand{\chaptitlefont}{\Huge\fontnaslova\bfseries}%
	\renewcommand{\printchaptertitle}[1]{%
		\chaptitlefont\hrule height 0.5pt \vspace{1em}%
		{##1}\vspace{1em}\hrule height 0.5pt%
	}%
	\renewcommand{\printchapternum}{%
		\chapnumfont\thechapter%
	}%
}

\pagestyle{ruled}
\chapterstyle{mychapterstyle}

\setsecheadstyle{\Large\fontnaslova\bfseries}
\setsubsecheadstyle{\large\fontnaslova\bfseries}
\setsubsubsecheadstyle{\normalsize\fontnaslova\bfseries}
% \paragraph kao run-in naslov: tekst se nastavlja u istom redu.
\setparaheadstyle{\normalsize\fontnaslova\bfseries}
\setparaindent{0pt}
\setbeforeparaskip{1.2ex plus 0.2ex minus 0.2ex}
\setafterparaskip{-0.5em}

% Broj stranice desno (oneside i neparne), lijevo na parnim (twoside).
\makeoddfoot{plain}{}{}{\thepage}
\makeevenfoot{plain}{\thepage}{}{}

% Numeracija i sadržaj
\setsecnumdepth{subsection}
\maxsecnumdepth{subsection}
\settocdepth{subsection}
\maxtocdepth{subsection}
